\section{Method: investment model and exact solvers}
\label{sec:method}

\paragraph{Setting.}
A leaderboard is a score matrix $S\in\mathbb{R}^{M\times n}$ ($M$ models, $n$ sub-datasets or ``tasks'', higher is better) and an aggregation rule $f$ that turns the scores of all models into a ranking. We fix a \emph{target} model $i$ and ask what the developer of $i$ must invest to make it strictly first under $f$. The other $M-1$ rows are fixed: rivals do not react, so the question is one-sided and not an equilibrium analysis.

\paragraph{Investment model.}
The developer may \emph{raise} the target's score on task $t$ to any $y_t\in[x_t,c_t]$, where $x=S_{i\cdot}$ and $c_t\ge x_t$ is a cap (a stand-in for ``training on data related to sub-dataset $t$''; scores are never lowered). A vector $y$ is \emph{successful} under $f$ if the target is strictly first with margin $\varepsilon$: for score rules, $f(y)\ge \max_{j\ne i}f(S_{j\cdot})+\varepsilon$ with $\varepsilon=10^{-3}$ points (a tie-breaking margin); the main experiments also report a statistical margin $k\cdot\mathrm{SE}$ ($k=1,2$) from binomial sampling of the per-task item counts, so that ``first'' is not an artefact of evaluation noise. We report three costs of $y$, never pooled into each other: $K=|\{t:y_t>x_t\}|$ (tasks touched; scale free), $L_1=\sum_t(y_t-x_t)$ (total points; for rule comparison also the ratio to the same target's cost under the mean rule), and $L_\infty=\max_t(y_t-x_t)$ (largest single-task jump, equal to the cost of a uniform improvement). If no $y$ within the caps succeeds, the outcome is \emph{infeasible}; it is reported as a rate and never replaced by a large number. The \emph{minimum investment} is the minimum of a cost over successful $y$.

\paragraph{Cap scenarios (the main modelling assumption).}
The mapping from training effort to score gain is unknown, so the cap is a scenario and results are never pooled across scenarios: \texttt{ceil} ($c_t=100$; isolates the rule), \texttt{front0}/\texttt{front2} (best existing score on $t$ plus $0$/$2$ points: ``cannot beat the state of the art by much''), \texttt{gain5/10/20} ($c_t=x_t+h$), and \texttt{emp95}/\texttt{emp99} (per-task 95th/99th percentile of observed fine-tune-minus-base gains on the Open LLM Leaderboard; applied to HELM only as a labelled transfer assumption, so HELM uses cap-free and \texttt{gain}-$h$ scenarios).

\paragraph{Aggregation rules.}
\emph{Score rules:} arithmetic mean (the official rule of the Open LLM Leaderboard; equal-weight on HELM subjects), item-count-weighted mean (HELM's official ``All Subjects''), power means with exponent $p<1$ (geometric $p\!\to\!0$, harmonic $p=-1$, interpolating to the minimum as $p\to-\infty$), median (even $n$: mean of the two middle values), trimmed mean, minimum. \emph{Rank rules:} Borda (mean rank), plurality (first-place count, ties split), and the Condorcet-winner requirement. The last is a \emph{labelled proxy} for pairwise rules: a Condorcet winner is first under every Condorcet-consistent rule (Copeland, Schulze, Kemeny), so its cost upper-bounds, but is not, their cost; we report how often the original leaderboard has no Condorcet winner.

\paragraph{Exact solvers.}
For each rule and cost we compute the exact minimum investment. Write $\tau=\max_{j\ne i}f(S_{j\cdot})$ and $b=\tau+\varepsilon$.

\begin{proposition}[mean, minimum, median]\label{prop:closed}
(i) \emph{Mean.} Let $\mathrm{need}=n b-\sum_t x_t$. If $\mathrm{need}\le 0$ the cost is $0$; otherwise first place is feasible iff $\sum_t(c_t-x_t)\ge\mathrm{need}$, the minimum $L_1$ is exactly $\mathrm{need}=n\cdot(\text{gap})+n\varepsilon$, and the minimum $K$ is the least $k$ such that the $k$ largest headrooms $c_t-x_t$ sum to at least $\mathrm{need}$ (so $K=1$ when one task has enough headroom).
(ii) \emph{Minimum.} First place is feasible iff $c_t\ge b$ for all $t$; then $L_1=\sum_t\max(0,b-x_t)$, $K=|\{t:x_t<b\}|$ and $L_\infty=b-\min_t x_t$.
(iii) \emph{Median, $n$ odd.} Let $m=(n+1)/2$. The median of $y$ is $\ge b$ iff at least $m$ coordinates are $\ge b$; hence $K=\max(0,m-|\{t:x_t\ge b\}|)$ and $L_1$ is the sum of the $m$ smallest gaps $\max(0,b-x_t)$ among tasks with $c_t\ge b$ (infeasible if fewer than $m$ such tasks).
\end{proposition}

Proofs are one line each: the mean constraint is $\sum_t y_t\ge nb$ and $L_1=\sum_t y_t-\sum_t x_t$; $\min_t y_t\ge b$ iff every $y_t\ge b$; the median is $\ge b$ iff at least $m$ order statistics are. Consequently the mean rule rewards \emph{concentration} (one large jump on the task with most headroom), whereas the median and the minimum force \emph{breadth}: a target below $b$ on every task needs $K\ge\lceil n/2\rceil$ under the median (for even $n$ the larger of the two middle values must reach $b$, so $K\ge n/2$; $K=n/2+1$ is not always needed) and $K=n$ under the minimum. For power means with $p<1$ the cost is minimised by equalising $y_t+s$ at a water level, found by bisection, and $K$ by a greedy rule on $g(c_t)-g(x_t)$. Rank rules have no closed form: Borda, plurality and the Condorcet requirement are solved exactly as mixed-integer programmes over candidate ``levels'' at which a raised score passes a rival on a task (rival constraints generated lazily; HiGHS through SciPy). Every solver is unit-tested against brute force on small random instances, and Proposition~\ref{prop:closed} is checked numerically against the solvers on every target of both datasets (Section~\ref{sec:exp}).

\paragraph{Targets and uncertainty.}
Targets are fixed ex ante: the challengers ranked 2--11 under the mean rule (primary), plus strata by initial rank band and gap. Uncertainty comes from bootstrap over targets, over tasks, and Monte-Carlo resampling of evaluation noise, as each claim demands; every aggregate is reported with its interval, its cap scenario and its feasibility rate.
